% Generated by scripts/assemble.py; corrections belong in source/corrections.json.
\examhead{178}{西北大学 2026 年数学分析试卷(回忆版)}{168}
% p169-b01; source PDF page 169
\blockanchor{p169-b01}
(注.题目来自考生回忆,部分表述或数据可能有误)\par
% p169-b02; source PDF page 169
\par\Needspace{5\baselineskip}
\blockanchor{p169-b02}
\majorquestion{一}(15 分) 解答如下问题:\par
% p169-b03; source PDF page 169
\blockanchor{p169-b03}
(1) 用极限的定义证明: 当 $\lim_{n\to \infty}a_n = 0$ 时, 有 $\lim_{n\to \infty}\frac{a_1 + a_2 + \cdots + a_n}{n} = 0.$\par
% p169-b04; source PDF page 169
\blockanchor{p169-b04}
(2) 若 $\lim_{n\to \infty}\frac{a_1 + a_2 + \cdots + a_n}{n} = 0,$ 是否有 $\lim_{n\to \infty}a_n = 0?$ 为什么？\par
% p169-b05; source PDF page 169
\blockanchor{p169-b05}
(3) 求 $\lim_{n\to \infty}\sqrt[n]{\frac{1}{n!}}.$\par
% p169-b06; source PDF page 169
\par\Needspace{5\baselineskip}
\blockanchor{p169-b06}
\majorquestion{二}(15 分) 已知 $f(x, y)$ 在 $(0, 0)$ 的某邻域连续, 且 $F(t) = \iint_{x^2 + y^2 \leq t^2} f(x, y) \, \mathrm{d}x \, \mathrm{d}y$, 求 $\lim_{t \to 0} \frac{F'(t)}{e^{3t} - 1}$.\par
% p169-b07; source PDF page 169
\par\Needspace{5\baselineskip}
\blockanchor{p169-b07}
\majorquestion{三}(15 分) 解答如下问题:\par
% p169-b08; source PDF page 169
\blockanchor{p169-b08}
(1) 证明: 若 $\lim_{n\to \infty}(a_{n + 1} - a_n) = a,$ 则 $\lim_{n\to \infty}\frac{a_n}{n} = a.$\par
% p169-b09; source PDF page 169
\blockanchor{p169-b09}
(2) 已知 $0 < x_{1} < 1$, $x_{n+1} = x_{n}(1 - x_{n})$, 证明: $\lim_{n \to \infty} nx_{n} = 1$.\par
% p169-b10; source PDF page 169
\par\Needspace{5\baselineskip}
\par\Needspace{6\baselineskip}
\blockanchor{p169-b10}
\majorquestion{四}(15 分) 设函数 $f(x)$ 在 $x = 0$ 的某邻域存在连续的一阶导数, 且 $f'(0) = 0$, $f''(0)$ 存在, 求极限\par
\nopagebreak[4]
% p169-b11; source PDF page 169
\blockanchor{p169-b11}
\displayformula{\lim _ {x \rightarrow 0} \frac {f (x) - f (\ln (1 + x))}{x ^ {3}}.}
% p169-b12; source PDF page 169
\par\Needspace{5\baselineskip}
\blockanchor{p169-b12}
\majorquestion{五}(15 分) 求 $\sum_{k=1}^{\infty}\left(\frac{(-1)^k}{2n+1}+\frac{1}{k(k+1)}\right)$.\par
% p169-b13; source PDF page 169
\par\Needspace{5\baselineskip}
\par\Needspace{6\baselineskip}
\blockanchor{p169-b13}
\majorquestion{六}(15 分) 求第二型曲面积分\par
\nopagebreak[4]
% p169-b14; source PDF page 169
\blockanchor{p169-b14}
\displayformula{\iint_ {S} (y + z) \mathrm{d} y \mathrm{d} z + (x ^ {2} + z ^ {2}) y \mathrm{d} z \mathrm{d} x + (x + y) \mathrm{d} x \mathrm{d} y.}
% p169-b15; source PDF page 169
\blockanchor{p169-b15}
\displayformula{\left(\text {另一个版本为} \iint_ {S} y z   \mathrm{d} y \mathrm{d} z + (x ^ {2} + z ^ {2}) y   \mathrm{d} z \mathrm{d} x + x y   \mathrm{d} x \mathrm{d} y\right)}
% p169-b16; source PDF page 169
\blockanchor{p169-b16}
其中 $S: 4 - y = x^{2} + z^{2} (y \geq 0)$ 的右侧.\par
% p169-b17; source PDF page 169
\par\Needspace{5\baselineskip}
\blockanchor{p169-b17}
\majorquestion{七}(15 分) 叙述并证明康托尔(Cantor)定理.\par
% p169-b18; source PDF page 169
\par\Needspace{5\baselineskip}
\blockanchor{p169-b18}
\majorquestion{八}(15 分) 设 $f(x, y) = |x - y| \varphi(x, y)$, 其中 $\varphi(x, y)$ 在 (0,0) 的某邻域内有定义, 且在 (0,0) 处连续. 证明: $\varphi(0,0) = 0$ 的充分必要条件是 $f(x, y)$ 在 (0,0) 处可微.\par
% p169-b19; source PDF page 169
\par\Needspace{5\baselineskip}
\blockanchor{p169-b19}
\majorquestion{九}（15分）已知黎曼函数\par
% p169-b20; source PDF page 169
\blockanchor{p169-b20}
$R(x)=\begin{cases}\frac{1}{q},&x=\frac{p}{q},\ \text{其中 }p,q\text{ 为正整数, 且 }(p,q)=1;\\0,&x\text{ 为无理数};\\1,&x=0,1.\end{cases}$\par (可能有误,黎曼函数在端点处取值应为0)\par
% p169-b21; source PDF page 169
\blockanchor{p169-b21}
讨论并说明  $R(x)$  在  $[0,1]$  上的可积性与连续性.\par
% p169-b22; source PDF page 169
\par\Needspace{5\baselineskip}
\blockanchor{p169-b22}
\majorquestion{十}(15 分) 已知 $g(0) = \frac{\sqrt{\pi}}{2}$, 求 $g(\alpha) = \int_{0}^{+\infty} e^{-x^2} \cos 2\alpha x \, \mathrm{d}x$.\par
